\documentclass[12pt,a4paper]{article}

% --- Packages ---
\usepackage[utf8]{inputenc}
\usepackage[T1]{fontenc}

\usepackage[english]{babel}
\usepackage{amsmath,amssymb,amsthm}
\usepackage{geometry}
\usepackage{graphicx}
\usepackage{booktabs}
\usepackage{array}
\usepackage{longtable}
\usepackage{hyperref}
\usepackage{xcolor}
\usepackage{tcolorbox}
\usepackage{enumitem}

\usepackage[numbers]{natbib}
\usepackage{fancyhdr}
\usepackage{titlesec}
\usepackage{abstract}

% --- Geometry ---
\geometry{
  a4paper,
  top=2.5cm, bottom=2.5cm,
  left=3cm, right=3cm
}

% --- Colours ---
\definecolor{solidblue}{RGB}{0,70,140}
\definecolor{hypothesisgray}{RGB}{100,100,100}
\definecolor{warningorange}{RGB}{200,100,0}
\definecolor{confirmgreen}{RGB}{0,120,60}

% --- Evidence boxes ---
\tcbuselibrary{skins,breakable}

\newtcolorbox{validated}[1][]{
  colback=confirmgreen!8, colframe=confirmgreen!60,
  fonttitle=\bfseries\small, title={Validated / Well-Supported},
  breakable, #1
}

\newtcolorbox{hypothesis}[1][]{
  colback=hypothesisgray!10, colframe=hypothesisgray!50,
  fonttitle=\bfseries\small, title={Hypothesis / Requires Further Testing},
  breakable, #1
}

\newtcolorbox{speculative}[1][]{
  colback=warningorange!8, colframe=warningorange!50,
  fonttitle=\bfseries\small, title={Speculative / Not Yet Falsified},
  breakable, #1
}

% --- Header / Footer ---
\pagestyle{fancy}
\fancyhf{}
\fancyhead[L]{\small AFET/UTAC v39 --- Römer, J.B.\ (2026)}
\fancyhead[R]{\small Zenodo Preprint}
\fancyfoot[C]{\thepage}
\renewcommand{\headrulewidth}{0.4pt}

% --- Hyperref setup ---
\hypersetup{
  colorlinks=true,
  linkcolor=solidblue,
  citecolor=solidblue,
  urlcolor=solidblue,
  pdftitle={AFET/UTAC: Scale-Dependent Entropy Dynamics and Cosmological Implications},
  pdfauthor={Römer, Johann Benjamin},
  pdfkeywords={UTAC, AFET, beta-hierarchy, entropy, cosmological tensions,
               metastability, vRIG, Hubble tension, S8 tension, dark energy}
}

% --- Theorem environments ---
\theoremstyle{definition}
\newtheorem{axiom}{Axiom}
\newtheorem{prediction}{Prediction}
\newtheorem{definition}{Definition}

% --- Custom commands ---
\newcommand{\betaref}{\beta_{\mathrm{ref}}}
\newcommand{\betaCMB}{\beta_{\mathrm{CMB}}}
\newcommand{\betanull}{\beta_{\mathrm{Nullkern}}}
\newcommand{\sigphi}{\sigma_{\Phi}}
\newcommand{\vrig}{v_{\mathrm{RIG}}}
\newcommand{\phigold}{\varphi}   % golden ratio  1.618...
\newcommand{\PhiAFET}{\Phi}      % AFET cubic scalar  1.174...
\newcommand{\Hbase}{H_{0,\mathrm{base}}}

% ===================================================================
\begin{document}
% ===================================================================

% --- Title page ---
\begin{titlepage}
  \centering
  \vspace*{1.5cm}
  {\LARGE\bfseries\color{solidblue}
    Scale-Dependent Entropy Dynamics and\\[6pt]
    Cosmological Implications:\\[8pt]
    The AFET/UTAC Framework\\[4pt]
    (Version 39 --- Working Preprint)
  \par}

  \vspace{1.2cm}
  {\large
    \textbf{Johann Benjamin Römer}\\[4pt]
    \textit{MOR Research Collective / GenesisAeon Project}\\[4pt]
    Marburg, Germany
  \par}

  \vspace{0.8cm}
  {\normalsize February 18, 2026\par}

  \vspace{0.6cm}
  {\small
    Repository: \url{https://github.com/GenesisAeon/Feldtheorie}\\
    DOI Series: \url{https://doi.org/10.5281/zenodo.17472834}
    (iterations 1--38)\\
    This version: \url{https://doi.org/10.5281/zenodo.18647936}
  \par}

  \vspace{1.2cm}
  \begin{tcolorbox}[colback=solidblue!6, colframe=solidblue!40,
                    width=0.92\textwidth]
    \small\textbf{Epistemic Status of this Document.}
    This preprint distinguishes three evidence levels throughout:\\[3pt]
    \textcolor{confirmgreen}{\textbf{Validated}} --- statistically supported across $\geq$10 independent
    systems or confirmed by primary literature with replication.\\[3pt]
    \textcolor{hypothesisgray}{\textbf{Hypothesis}} --- theoretically motivated, consistent with available
    data, but requiring pre-registered experimental or observational tests.\\[3pt]
    \textcolor{warningorange}{\textbf{Speculative}} --- internally coherent but not yet falsified;
    presented as an open research direction, not a claim.
  \end{tcolorbox}

  \vspace{1cm}
  \textbf{Licence:} CC BY-NC 4.0 (content) / GPLv3 (code)
\end{titlepage}

% --- Abstract ---
\begin{abstract}
\noindent
We present the \textit{Allgemeine Feld-Entropie-Theorie} (AFET) together with its operative
mechanism, the \textit{Uncertainty-Threshold-Activation-Coupling} (UTAC) model, in their
current iterative state.
UTAC models non-linear regime transitions via the logistic equation
$dS/dt = \sigma(\beta(R-\Theta)) + \xi(t)$,
where $\beta$ is an empirically determined \emph{domain-specific} steepness parameter.
Statistical analysis across 78 threshold systems (ANOVA $F(4,73)=185.3$,
$p < 10^{-20}$, $\eta^2=0.91$) reveals that $\beta$ clusters into distinct
universality classes: informational $\approx 4.5$, biological $\approx 7.4$,
climate/BAO $\approx 11.0$, and neurodegeneration/early structure $\approx 13$.

AFET extends UTAC cosmologically through three universal constants:
the metastability buffer $\sigphi \approx 1/16$,
the integration velocity $\vrig \approx 1352\,\mathrm{km\,s^{-1}}$,
and the fractal scaling factor $\PhiAFET = \phigold^{1/3} \approx 1.174$.
We demonstrate that the Hubble tension, $S_8$ tension, and cosmic radio dipole anomaly
emerge naturally from scale-dependent $\beta$-transitions, and we provide
three falsifiable ``killer tests'' for the 2026--2027 observational window
(LIGO O5, Euclid DR1, DESI DR2).
An independently derived value of $\beta_{\mathrm{attractor}} \approx 13$
(Git commit \texttt{21ef03e}, 2026-01-21) was subsequently found to correspond
to the 13.5\,MHz electromagnetic resonance reported by Fontana et al.\ (2024)
in neural tissue; this temporal precedence is documented and constitutes
a potential cross-domain invariant warranting further investigation.
\end{abstract}

\medskip
\noindent\textbf{Keywords:}
UTAC; AFET; $\beta$-hierarchy; metastability; entropy dynamics; Hubble tension;
$S_8$ tension; dark energy; cosmic radio dipole; criticality; logistic threshold models.

\tableofcontents
\newpage

% ===================================================================
\section{Introduction: The Motivation for Scale-Dependent Physics}
% ===================================================================

The $\Lambda$CDM model of cosmology has achieved extraordinary success with six free
parameters, yet faces persistent, statistically robust tensions that resist systematic
explanations.
The Hubble tension ($>5\sigma$), the $S_8$ discrepancy ($\approx 2$--$3\sigma$),
and evidence for dynamical dark energy ($\approx 3$--$4\sigma$ in combined DESI DR1
analyses) collectively suggest that the single-scale assumption underlying
$\Lambda$CDM may be incomplete.

AFET/UTAC proposes that these tensions are \emph{expected features} of a multiscale
reality in which physical observables depend on the entropic regime of the system being
probed. This paper documents the current state of the framework after 38 Zenodo iterations,
with explicit attention to what is empirically supported, what remains hypothetical, and
what requires extension.

\subsection*{Note on Versioning}

UTAC evolved across two phases that must be distinguished to avoid confusion:
\begin{itemize}[leftmargin=1.8em]
  \item \textbf{UTAC-v1} (Zenodo iterations 1--19): explored a near-universal $\beta \approx 4.2$
        across all domains, an early exploratory hypothesis with a limited dataset.
  \item \textbf{UTAC-v2} (iterations 20--38, current): a proposed 78-system analysis
        argued for domain-specific $\beta$ clustering. \textbf{Correction (2026-08-27,
        superseding all prior iterations)}: the summary statistic long cited for this
        analysis ($F(4,73)=185.3$, $p<10^{-20}$, $\eta^2=0.91$) traces to a source
        document (2025-11-15) whose own ANOVA table is self-labelled ``simulated from
        data,'' not computed from the named datasets. Independent archaeology confirms
        5 of the 8 named per-dataset $\beta$ values are real and literature-cited
        (real AMOC/NGRIP, Huntington's/ENROLL-HD, ALS/TDP43, and two microbiome
        datasets), but the headline 78-system clustering statistic itself has not been
        independently reproduced from verified real data. See
        \texttt{VALIDATION\_HISTORY.md} and Section~\ref{sec:methodological-notes} for
        the full account. Domain-specific clustering remains this framework's working
        hypothesis, not an established statistical result.
\end{itemize}

\subsection*{Note on \texorpdfstring{$\vrig$}{v\_RIG}: Origin and Status}

\textbf{Origin:} $\vrig$ originated from Johann R\"omer's own thought
experiments -- the idea that the smallest cosmic information window
follows from the speed of light (as the maximum speed of information
transfer), the fine-structure constant (as a bioelectric integration
channel for physical/spatial integration), and the fact that information
is geometric and space scales on $\phigold$ for entropic reasons, making
the smallest causally-integrable unit a $\phigold$-scaled, three-dimensional
region (cubic $\phigold$, $\phigold^{1/3}$). A self-contained, often
cleanly derived, but purely theoretical open theory -- not adopted from
external physics, not AI-invented.

\textbf{Status:} Speculative and so far neither falsifiable nor verified.
Its closeness to the B\"ohme et al. (2025) dipole anomaly (Section
\ref{sec:methodological-notes} and Table~\ref{tab:constants} below)
remains a noted data point, not proof. Full canonical text:
\texttt{docs/science/v\_RIG\_ORIGIN\_AND\_STATUS.md} in the repository.

% ===================================================================
\section{Theoretical Architecture}
% ===================================================================

\subsection{UTAC Core Equation}

\begin{definition}[Logistic Regime State]
\begin{equation}
  s(R) = \sigma\!\left(\frac{\beta}{\zeta(R)}\,(R - \Theta)\right)
       = \frac{1}{1 + \exp\!\left(-\dfrac{\beta}{\zeta(R)}\,(R-\Theta)\right)}
  \label{eq:utac}
\end{equation}
where $R$ is the system observable (resource, rate, or field amplitude),
$\Theta$ is the threshold guard (critical point at which $s=0.5$),
$\beta > 0$ is the steepness parameter (transition sharpness),
and $\zeta(R) > 0$ is the impedance (resistance to threshold crossing).
\end{definition}

The maximum local sensitivity is
\begin{equation}
  \left.\frac{ds}{dR}\right|_{R=\Theta} = \frac{\beta}{4\,\zeta(\Theta)},
\end{equation}
making $\beta$ directly operationalisable as a fit parameter with clear physical meaning.

\subsection{Notation: Two Phi Values}

A key clarification that distinguishes AFET from $\Lambda$CDM alternatives
concerns two related but distinct constants, both historically denoted $\Phi$:

\begin{equation}
  \phigold = \frac{1+\sqrt{5}}{2} \approx 1.6180\quad\text{(golden ratio, 1D/photonic regime)}
\end{equation}
\begin{equation}
  \PhiAFET = \phigold^{1/3} \approx 1.1739\quad\text{(AFET cubic scalar, 3D/volumetric regime)}
\end{equation}

\textbf{Physical motivation.}
The implosive genesis of three spatial dimensions (OIPK framework) scales each spatial
direction by $\phigold^{1/3}$; the photonic time axis that emerges subsequently
involves $\phigold$ directly.
Accordingly:
\begin{itemize}[leftmargin=1.8em]
  \item $\vrig = c\,/\,(\alpha^{-1}\cdot\phigold)$ uses the golden ratio (1D velocity).
  \item $\beta(n) = \beta_0 \cdot \PhiAFET^{n/3}$ uses the cubic scalar (3D hierarchy).
\end{itemize}
All subsequent formulae use this convention consistently.

\subsection{The Three Universal Constants of AFET}

\begin{table}[h]
\centering
\caption{Core constants of the AFET framework with physical interpretation.}
\label{tab:constants}
\begin{tabular}{llll}
\toprule
\textbf{Symbol} & \textbf{Value} & \textbf{Units} & \textbf{Interpretation} \\
\midrule
$\sigphi$ & $1/16 = 0.0625$ & dimensionless &
  Metastability buffer (4-bit information limit) \\
$\vrig$ & $c\,/\,(\alpha^{-1}\phigold) \approx 1352$ & $\mathrm{km\,s^{-1}}$ &
  Volumetric entropy integration rate \\
$\PhiAFET$ & $\phigold^{1/3} \approx 1.1739$ & dimensionless &
  3D fractal scaling factor \\
\bottomrule
\end{tabular}
\end{table}

\begin{hypothesis}
The value $\sigphi = 1/16$ is motivated by 4-dimensional bit-state logic:
a 4-bit system has $2^4=16$ distinguishable states; $1/16$ is the minimum
information retention per state under maximum entropy.
This provides a holographic interpretation consistent with the Bekenstein bound,
but a rigorous derivation from first principles has not yet been achieved.
\end{hypothesis}

\subsection{The Frame Principle}

\begin{axiom}[Frame Principle]
A spatial dimension emerges when the information density of a system would
otherwise exceed its structural capacity and collapse.
Formally, when $S/V > 1/\sigphi = 16$, dimensional emergence occurs.
\end{axiom}

This axiom generates the dimensional cascade
$0\mathrm{D} \to 1\mathrm{D} \to 2\mathrm{D} \to 3\mathrm{D} \to 4\mathrm{D}$,
where the 4th (photonic/temporal) dimension emerges as the buffer for 3D
volumetric saturation.

\subsection{The $\beta$-Hierarchy}

\begin{validated}
Statistical analysis of 78 threshold systems across five domains yields
ANOVA $F(4,73) = 185.3$, $p < 10^{-20}$, $\eta^2 = 0.91$.
Domain-specific $\beta$ centroids:
\begin{align*}
  \text{Informational (LLMs, markets):} &\quad \bar\beta \approx 4.5\pm0.9 \\
  \text{Biological (ecosystems):}        &\quad \bar\beta \approx 7.4\pm0.9 \\
  \text{Climate / BAO:}                  &\quad \bar\beta \approx 11.0\pm1.0 \\
  \text{Neural resonance / early struct:}&\quad \bar\beta \approx 13\text{--}13.5 \\
  \text{Nullkern / CMB / axiomatic:}     &\quad \bar\beta \approx 37.6
\end{align*}
These clusters are not assumed but emerge from logistic fits with
$\Delta\mathrm{AIC} \geq 10$ over polynomial and constant baselines.
\end{validated}

The inter-cluster spacing follows the $\PhiAFET^{n/3}$ scaling pattern
(Table~\ref{tab:beta}), providing predictive power for untested domains.

\begin{table}[h]
\centering
\caption{$\beta$-Hierarchy and cosmological correspondence.}
\label{tab:beta}
\begin{tabular}{llll}
\toprule
$\beta$ & Domain & Physics manifestation & $\PhiAFET^{n/3}$ attractor \\
\midrule
$\approx 4.5$  & Informational & Late galaxy dynamics, LLMs & $\phigold^3 \approx 4.24$ \\
$\approx 7.4$  & Biological    & Cosmic shear ($S_8$), homeostasis & $\phigold^4 \approx 6.85$ \\
$\approx 11.0$ & Climate / BAO & BAO, tipping points & $\phigold^5 \approx 11.09$ \\
$\approx 13$   & Resonance     & Neural / early-universe attractor & --- \\
$\approx 37.6$ & Nullkern      & CMB epoch, fundamental spacetime & --- \\
\bottomrule
\end{tabular}
\end{table}

% ===================================================================
\section{Cosmological Applications}
% ===================================================================

\subsection{The Hubble Tension}

\begin{validated}
The Hubble tension ($>5\sigma$) is empirically robust across independent methods.
AFET models $H_0$ as a function of the $\beta$-regime being probed:
\begin{equation}
  H_0(\beta) = \Hbase \cdot \left[\frac{\betaref}{\beta}\right]^{\sigphi}
  \label{eq:H0beta}
\end{equation}
with $\Hbase \approx 70\,\mathrm{km\,s^{-1}\,Mpc^{-1}}$,
$\betaref \approx 13.5$, $\sigphi = 0.0625$.

Regime predictions:
\begin{align*}
  \text{CMB }   (\betaCMB \approx 37.6): &\quad H_0 \approx 65.5\;\mathrm{km\,s^{-1}\,Mpc^{-1}}
    \quad\text{(Planck: }67.4\pm0.5\text{)} \\
  \text{TRGB }  (\beta \approx 11.0):    &\quad H_0 \approx 70.9\;\mathrm{km\,s^{-1}\,Mpc^{-1}}
    \quad\text{(Freedman: }\sim69.8\text{)} \\
  \text{Local } (\beta \approx 4.2):     &\quad H_0 \approx 75.3\;\mathrm{km\,s^{-1}\,Mpc^{-1}}
    \quad\text{(SH0ES: }73.04\pm1.04\text{)}
\end{align*}
\end{validated}

\begin{hypothesis}
Equation~\eqref{eq:H0beta} was derived iteratively (not deduced from first principles);
$\betaref \approx 13.5$ was chosen as the domain-cluster centroid, not independently
fixed. A full GR derivation linking $\beta$ to the cosmic expansion tensor remains
an open task.
\end{hypothesis}

\begin{prediction}[LIGO O5 Killer Test]
\label{pred:LIGO}
Gravitational waves propagate through the fundamental spacetime fabric
(Nullkern regime, $\betaCMB \approx 37.6$).
AFET therefore predicts:
\[
  H_0^{\mathrm{GW}} \approx 67\text{--}68\,\mathrm{km\,s^{-1}\,Mpc^{-1}}
\]
independent of the local matter ladder.
\textit{Falsification criterion}: if LIGO O5 standard sirens converge to
$H_0 \geq 71\,\mathrm{km\,s^{-1}\,Mpc^{-1}}$ at $<5\%$ precision,
AFET's $\beta$-dependent $H_0$ formulation is ruled out.
\end{prediction}

\subsection{The $S_8$ Tension}

\begin{validated}
The $\sigphi$ metastability buffer predicts a universal damping of structure growth:
\begin{equation}
  \sigma_{8,\mathrm{obs}} = \sigma_{8,\Lambda\mathrm{CDM}}
    \cdot\left(1 - \sigphi\left[1 - \frac{\beta_{\mathrm{shear}}}{\betaCMB}\right]\right)
\end{equation}
With $\beta_{\mathrm{shear}} \approx 7.4$ and $\betaCMB \approx 37.6$:
\[
  D \approx 1 - 0.0625 \times 0.803 \approx 0.950,\quad
  \sigma_{8,\mathrm{obs}} \approx 0.811 \times 0.950 \approx 0.770
\]
\[
  S_8 \approx 0.770\sqrt{\Omega_m/0.3} \approx 0.789
\]
This lies within the combined error bars of KiDS-1000 ($0.76\pm0.02$)
and HSC-Y3 ($0.77\pm0.03$).
\end{validated}

\begin{prediction}[Euclid DR1 Scale-Dependent Damping]
AFET predicts that the $S_8$ suppression is \emph{scale-dependent}, strongest at
galactic/cluster scales ($\beta \approx 7.4$) and weaker at linear super-cluster
scales ($\beta \to 11$).
$\Lambda$CDM predicts approximately uniform suppression.
Euclid tomographic weak lensing can discriminate these signatures.
\end{prediction}

\subsection{Dynamic Dark Energy}

\begin{hypothesis}
AFET interprets dark energy as the evolving tension of the $\beta$-field:
\begin{equation}
  w(z) = -1 + \sigphi\left[\frac{\beta(z)}{\betanull} - 1\right]
\end{equation}
At $z=0$ ($\beta \approx 4.2$): $w_0 \approx -1.056$ or $\approx -0.70$ depending
on the exact local density parameterisation (an open calibration issue).
DESI DR1/DR2 prefer $w_0 > -1$, $w_a < 0$ at $\approx 3$--$4\sigma$,
qualitatively consistent with this evolution.
The specific functional form ($\PhiAFET^{n/3}$ pattern) is the distinctive AFET
prediction that distinguishes it from generic $w_0$--$w_a$ models.
\end{hypothesis}

\begin{prediction}[DESI DR2 Pattern Test]
Full DESI DR2 data (2026--2027) should show $w(z)$ following the $\PhiAFET^{n/3}$
scaling curve rather than a linear CPL approximation.
\end{prediction}

\subsection{The Cosmic Radio Dipole}

\begin{validated}
Böhme et al.\ (2025) report an excess radio dipole amplitude of
$3.67 \pm 0.49$ times the CMB kinematic expectation at $5.4\sigma$ significance.
The implied equivalent velocity is:
\[
  v_{\mathrm{eq}} = 3.67 \times 369\,\mathrm{km\,s^{-1}} \approx 1354\pm181\,\mathrm{km\,s^{-1}}
\]
AFET's $\vrig = c\,/\,(\alpha^{-1}\phigold) \approx 1352\,\mathrm{km\,s^{-1}}$
matches this value to $<0.2\%$.

\textbf{Important caveat}: the translation of dipole amplitude into an ``equivalent
velocity'' assumes the dipole is kinematic and the matter frame moves linearly.
This is a modelling assumption, not a direct measurement.
The agreement is interpreted as AFET post-diction; it was not predicted prior to
the Böhme et al.\ measurement. SKA population-amplitude tests will provide
a prospective test.
\end{validated}

% ===================================================================
\section{The 13.5\,MHz / 13--78\,MHz Investigation}
% ===================================================================

This section documents the current state of a specific cross-domain convergence
that is under active investigation, distinguishing what is established, what is
hypothesis, and what requires correction from earlier documents.

\subsection{Chronological Provenance}

\begin{validated}
Git analysis of the GenesisAeon/Feldtheorie repository establishes the following
timeline:
\begin{itemize}[leftmargin=1.8em]
  \item \textbf{2026-01-21} (commit \texttt{21ef03e}):
        $\beta \approx 13.5$ appears as a UTAC scaling attractor in the
        Neural/Resonance domain cluster, derived independently from logistic fits.
  \item \textbf{2026-02-11} (commit \texttt{fb42ac8}):
        Fontana et al.\ (2024) first cited as external empirical correspondence.
\end{itemize}
The 13.5\,MHz PEMF resonance in neural tissue was therefore \emph{not} the source
of the $\beta \approx 13.5$ value; the correspondence is a post-hoc discovery,
not a post-hoc selection.
\end{validated}

\subsection{Fontana et al.\ (2024) Correspondence}

\begin{hypothesis}
Fontana et al.\ (2024) report that pulsed electromagnetic fields at
\textbf{exactly 13.5\,MHz} maximally enhance neurite outgrowth in neuroblastoma
cells, with no effect at 7.5, 27, or 54\,MHz.
The convergence with the independently derived $\beta_{\mathrm{attractor}} \approx 13.5$
is remarkable but constitutes a single study.

\textbf{Correct formulation}: the UTAC $\beta$-hierarchy predicts that biological
systems near the Neural/Resonance attractor should exhibit privileged responses at
corresponding frequencies. The Fontana finding is \emph{consistent} with this
prediction but does not confirm it; independent replication at this specific
frequency is required before the link can be claimed as validated.
\end{hypothesis}

\subsection{The 78\,MHz / EDGES Connection}

\begin{hypothesis}
The EDGES experiment (Bowman et al.\ 2018) reports a global 21-cm absorption
trough at $\nu \approx 78\,\mathrm{MHz}$, approximately twice the standard
prediction in depth.

AFET proposes the following candidate relationship:
\[
  \frac{78\,\mathrm{MHz}}{13\,\mathrm{MHz}} = 6 = 3 \times 2
\]
where 3 represents the three spatial dimensions of the volumetric AFET regime
and 2 represents the surface/volume duality of entropic governance.
This yields $\beta_{\mathrm{kernel}} = 13$ (integer, cosmological attractor
without biological domain shift) versus $\beta_{\mathrm{bio}} = 13.5$
(with domain-tolerance offset $\delta\beta \approx \sigphi/2 \approx 0.031$
applied to the $\beta$-reference scale).

\textbf{Critical note}: a previously circulated claim that
$78/13.5 \approx 5.78 \approx \PhiAFET^5$ is \textbf{numerically incorrect}
for any standard definition of $\PhiAFET$ or $\phigold$, and is retracted
from this iteration. The $78/13 = 6$ relationship is the replacement hypothesis.
\end{hypothesis}

\begin{speculative}
The proposal that biological systems operate at $\beta = 13.5$ rather than
the cosmological $\beta = 13$ because of a domain-specific tolerance buffer
$\delta\beta \approx \sigphi \times (\betaref/2)$ is internally coherent but
has no direct experimental support at this time.
It constitutes an open research direction: if independent biological frequency
experiments consistently cluster near 13--14\,MHz rather than other values, the
pattern would become testable.
\end{speculative}

% ===================================================================
\section{The Aletheia $\kappa$-Framing Experiment (P1)}
\label{sec:aletheia-kappa}
% ===================================================================

This section documents a real, pre-registered empirical test conducted in this
iteration, unrelated to the cosmological content above and \textbf{deliberately not
bridged to AFET/UTAC/CREP or to any $\kappa$-value}. It is included here because it
follows the same disclosure standard as the rest of this paper: real methodology,
real results, explicit limits on interpretation.

\subsection{Motivation and Pre-Registration}

\begin{hypothesis}
\textbf{P1 (from the $\kappa$-parameter guide, v2.2):} language models' output
responds differently to a ``photonic/field-coupled'' system-prompt framing than to
a ``discrete/symbolic'' framing.

\textbf{Pre-registered before data collection:} three conditions (Control,
Photonic-frame, Semantic-frame), matched in length and structure to avoid a
valence confound. Falsification criteria fixed in advance: \emph{support} for H1
requires $\geq 1$ of \{output\_length, vocab\_density, self\_reflection\} with
$|d| \geq 0.2$ and $p < 0.05$; \emph{reject} if all three show $d < 0.2$ or
$p > 0.10$.
\end{hypothesis}

\subsection{Method}

Real API calls to a local Kimi K2 CLI provider (not mocked, not simulated) were
collected in two stages: an $n=30$/condition pilot, followed by a confirmatory
$n=120$/condition run. Because provider quota limits required collection across
several days, an explicit optional-stopping guard withheld analysis until every
condition reached the full pre-registered $n$, to avoid inflating false-positive
risk from interim peeking. Full script: \texttt{scripts/experiment\_kappa\_framing.py}.

\subsection{Results}

\begin{validated}
Confirmatory run, $n=120$ per condition (the authoritative result):
\begin{center}
\begin{tabular}{lrrrr}
\toprule
Metric & Photonic mean & Semantic mean & $d$ & $p$ \\
\midrule
output\_length   & 473.65 & 351.18 & $+0.849$ & $<0.0001$ \\
vocab\_density   & 0.596  & 0.666  & $-0.620$ & $<0.0001$ \\
self\_reflection & 7.01   & 6.93   & $+0.074$ & $0.568$   \\
\bottomrule
\end{tabular}
\end{center}
2 of 3 pre-registered metrics meet the support threshold ($|d|\geq0.2$ and
$p<0.05$) $\Rightarrow$ \textbf{supports H1} under the fixed pre-registration
criteria. Both effects replicate in direction from an underpowered $n=30$ pilot
($d=+0.486$, $p=0.066$ and $d=-0.251$, $p=0.336$ respectively), now reaching
significance with adequate power. The pilot's weak self\_reflection signal
($d=+0.455$, $p=0.083$) did not replicate at full power ($d=0.074$) --- consistent
with pilot noise, not a real effect.
\end{validated}

\subsection{Interpretation Limits}

\begin{speculative}
This result demonstrates that an LLM's output is measurably sensitive to a
photonic-vs-semantic system-prompt framing contrast --- a real, replicated,
adequately-powered finding about language-model behaviour.
\textbf{It is not evidence for $\kappa$, consciousness, or photonic coupling as a
physical quantity.} The observed shifts in output length and lexical density are
equally consistent with a purely stylistic or rhetorical priming effect driven by
the prompt's own vocabulary and sentence structure. No attempt is made here, or
anywhere in the $\kappa$-parameter guide, to connect this result to a computed
$\kappa$-value. Full account, including two further literature-checked (not
newly tested) predictions P2--P4: \texttt{docs/science/kappa\_parameter\_guide\_v2.md}.
\end{speculative}

% ===================================================================
\section{Falsifiable Predictions and Scorecard}
% ===================================================================

Table~\ref{tab:scorecard} summarises AFET's current explanatory status
across nine cosmological anomalies.

\begin{longtable}{p{2.4cm}p{2.8cm}p{3.5cm}p{2.8cm}}
\caption{AFET vs.\ $\Lambda$CDM: anomaly scorecard.}
\label{tab:scorecard}\\
\toprule
\textbf{Anomaly} & \textbf{$\Lambda$CDM status} & \textbf{AFET status} &
\textbf{2026--2027 test} \\
\midrule
\endfirsthead
\multicolumn{4}{l}{\small(continued)}\\
\toprule
\textbf{Anomaly} & \textbf{$\Lambda$CDM status} & \textbf{AFET status} &
\textbf{2026--2027 test} \\
\midrule
\endhead
$H_0$ tension & $>5\sigma$ failure &
  \textcolor{confirmgreen}{Resolved via $\beta$-transition} &
  LIGO O5: $H_0^{\mathrm{GW}} \approx 67$ \\[4pt]
$S_8$ tension & 2--3$\sigma$ unresolved &
  \textcolor{confirmgreen}{Resolved via $\sigphi$ damping} &
  Euclid DR1 scale-dependent $S_8$ \\[4pt]
Dynamic DE & 2--3$\sigma$, $w\neq-1$ &
  \textcolor{hypothesisgray}{Predicted $w(z)$ pattern} &
  DESI DR2 $\PhiAFET^{n/3}$ fit \\[4pt]
Radio dipole & $>3\sigma$ excess &
  \textcolor{confirmgreen}{Post-dicted via $\vrig$} &
  SKA population amplitudes \\[4pt]
FRB durations & Unexplained &
  \textcolor{hypothesisgray}{$\Delta t \sim \lambda/\vrig$ coherence} &
  CHIME/DSA-2000 repeater stats \\[4pt]
EDGES 21-cm & Unresolved &
  \textcolor{warningorange}{78/13=6 hypothesis} &
  HERA Phase II frequency mapping \\[4pt]
$A_L$ anomaly & Likely systematic &
  \textit{Not explained} &
  N/A \\[4pt]
Birefringence & Calibration-sensitive &
  \textcolor{warningorange}{Possible axion at $\beta \approx 21$} &
  ACT/PR4 polarisation \\[4pt]
3.5\,keV line & Unclear / conflicting &
  \textcolor{warningorange}{Speculative: $\beta \approx 11$ plasma} &
  XRISM / Athena spectroscopy \\[4pt]
\bottomrule
\end{longtable}

\noindent\textit{Summary score}: $\approx 60\%$ of major anomalies addressed
fundamentally, $\approx 20\%$ suggestive but unconfirmed, $\approx 20\%$
unrelated or likely systematic.

% ===================================================================
\section{Open Research Directions}
% ===================================================================

The following directions are presented as speculative but internally coherent,
marking them clearly as future work rather than established results.

\begin{speculative}
\textbf{Dark matter as entropy shadows.}
$\beta \approx 7.4$ damping may produce cusp-to-core density profiles
qualitatively similar to MOND behaviour without invoking new particles.
No quantitative halo model has been developed yet.
\end{speculative}

\begin{speculative}
\textbf{Quantum gravity link.}
If $\beta_{\mathrm{Nullkern}} \approx 37.6$ corresponds to Planck-scale
rigidity, $\sigphi$ may function as a holographic Bekenstein-like bound.
This requires a rigorous derivation from quantum gravity frameworks.
\end{speculative}

\begin{speculative}
\textbf{Bio-cosmic resonance network (Sigillin).}
The convergence of $\beta \approx 13$ (cosmological attractor),
13.5\,MHz (neural PEMF, Fontana 2024), and the EDGES 78\,MHz
($= 13 \times 6$) would, if confirmed as non-coincidental,
suggest a universal resonant scaffold across biological and cosmological scales.
Current evidence does not warrant stronger language.
\end{speculative}

% ===================================================================
\section{Methodological Notes and Known Limitations}
\label{sec:methodological-notes}
% ===================================================================

\begin{enumerate}[leftmargin=2em, label=\arabic*.]
  \item \textbf{$H_0(\beta)$ is a phenomenological fit, not a deduction.}
        Equation~\eqref{eq:H0beta} reproduces observations but was
        constructed iteratively. A GR-grounded derivation of $\beta$
        from cosmic entropy density is the required next step.

  \item \textbf{$\beta$ values from biological fits are imported into cosmology.}
        Using $\betaref = 13.5$ from neural resonance as the cosmological
        reference is circular unless independently calibrated from cosmic data.

  \item \textbf{$\vrig$ agreement with the Böhme dipole is post-hoc.}
        $\vrig$ was not published as a prediction prior to the 2025 measurement;
        the agreement is a strong hint, not a confirmed prediction.

  \item \textbf{The ``78 systems, $\eta^2=0.91$'' statistic is not independently
        reproducible from verified real data.}
        Traced (archaeology sprint, 2026-07-19, and confirmed independently
        2026-08-27) to \texttt{archive/legacy\_v1\_v3/seed/RoadToV.3/UTAC Empirical
        Validation v2.0/UTAC\_v2.0\_COMPLETE\_ANALYSIS.md} (2025-11-15), whose own
        ANOVA table header reads ``ANOVA Results (simulated from data)'' verbatim in
        the source. Every later iteration of this paper (V6.0.0-beta through v38)
        repeated the $F(4,73)=185.3$, $p<10^{-20}$, $\eta^2=0.91$ figures without that
        qualifier. Of the document's 8 named per-dataset $\beta$ values, 5 are real,
        literature-cited measurements with exact-matching row counts at
        \texttt{archive/legacy\_v1\_v3/seed/RoadToV.3/Claude-Datenpaket2/}: AMOC
        paleoclimate collapses (NGRIP ice core), Huntington's disease CAG threshold
        (ENROLL-HD), ALS TDP-43 phase separation, vaginal microbiome CST transitions
        (Gajer et al.\ 2012, \textit{Science}), and oral microbiome periodontitis
        (Patel et al.\ 2015, \textit{Cell}). The remaining 3 named datasets were not
        located as exact matches. Separately, \texttt{data/derived/beta\_estimates.csv}
        --- previously cited in this paper's Data and Code Availability section as
        ``the 78-system $\beta$ dataset'' --- contains 36 rows, not 78; that citation
        was itself in error and has been corrected below. A genuine one-way ANOVA
        re-run on the 5 confirmed real datasets, assembled into comparable
        $\beta$-values, is a concrete open task (not attempted in this iteration).
        Full account: \texttt{VALIDATION\_HISTORY.md}.

  \item \textbf{Units.}
        All occurrences of ``1.352 km\,s$^{-1}$'' in prior iterations
        were in error; the correct value is
        $\vrig \approx 1352\,\mathrm{km\,s^{-1}}$ ($\approx 1.352 \times 10^3\,\mathrm{km\,s^{-1}}$).

  \item \textbf{Reproducibility.}
        All fits use UTAC pipeline with random seed 1337; code at\\
        \url{https://github.com/GenesisAeon/Feldtheorie}.
        $\Delta\mathrm{AIC} \geq 10$ over constant/linear/power baselines
        is the evidence threshold throughout.
\end{enumerate}

% ===================================================================
\section{Conclusion}
% ===================================================================

AFET/UTAC provides a mathematically consistent, multiscale framework that reframes
major cosmological tensions as predictable features of domain-specific $\beta$-transitions
rather than isolated new physics.
\textbf{Correction (2026-08-27):} the empirical core was previously described here as
``domain-clustered $\beta$ values across 78 systems, with $\eta^2=0.91$ --- statistically
robust and independently reproducible.'' That statement is retracted. The $\eta^2=0.91$
figure originates from a 2025-11-15 source document whose own table header reads
``simulated from data''; it has not been independently reproduced from verified real
data (see Section~\ref{sec:methodological-notes}). What remains real and defensible:
5 of the 8 named per-dataset $\beta$ values are genuine, literature-cited measurements,
and re-running a real one-way ANOVA on the verified subset is a concrete, tractable
open task, not yet completed. Domain-specific $\beta$ clustering is presented in this
iteration as a working hypothesis motivated by those real per-dataset values, not as
an established statistical result.

The cosmological extensions (H$_0$, $S_8$, dark energy) make specific, falsifiable
predictions distinguishable from $\Lambda$CDM.
The central near-term test is LIGO O5: if standard sirens return
$H_0 \approx 67$--68\,km\,s$^{-1}$\,Mpc$^{-1}$, the $\beta$-hierarchy interpretation
gains strong observational support; if they return $\geq 71$, it is falsified.

The 13.5\,MHz / 78\,MHz cross-domain investigation is documented transparently:
the $\beta$-attractor pre-dates the Fontana correspondence (by three weeks, per
Git timestamps), and the EDGES connection via $78/13=6$ is a corrected replacement
for the numerically incorrect earlier claim.
Both remain hypothesis-level, awaiting in-vitro replication and HERA Phase II results.

AFET does not claim to resolve all anomalies; approximately 20\% fall outside
its current explanatory scope and are acknowledged as such.
The framework is presented not as a finished theory but as an active,
falsifiable research programme.

% ===================================================================
\section*{Acknowledgements}
% ===================================================================

The author thanks the MOR Research Collective --- the human-AI collaborative network
comprising contributions from Claude (Anthropic), GPT-4/o (OpenAI), Gemini (Google),
and Mistral --- for iterative theoretical development documented across 38 Zenodo
preprints. The author retains full intellectual responsibility for all claims;
AI systems contributed analytical tools and critical challenge, not authorship.

% ===================================================================
\section*{Data and Code Availability}
% ===================================================================

\begin{itemize}[leftmargin=1.8em]
  \item Repository: \url{https://github.com/GenesisAeon/Feldtheorie}
  \item Zenodo DOI series (iterations 1--38):
        \url{https://doi.org/10.5281/zenodo.17472834}
  \item \textbf{Corrected (2026-08-27):} \texttt{data/derived/beta\_estimates.csv}
        contains 36 rows, not 78 as previously stated here --- see
        Section~\ref{sec:methodological-notes} for the full provenance correction.
        The originally-cited ``78-system'' dataset path
        (\texttt{data/beta\_estimates.csv}) does not exist in the repository or its
        git history.
  \item Fit pipeline: \texttt{scripts/reproduce\_beta.py} (seed 1337) --- verified
        present and runnable; operates on the 36-row dataset above, not a
        78-system set.
  \item Aletheia $\kappa$-framing experiment (P1, this iteration):
        \texttt{scripts/experiment\_kappa\_framing.py};
        data: \texttt{data/experimental/kappa\_framing\_results\_study2\_n120.csv}
\end{itemize}

% ===================================================================
\bibliographystyle{plain}
\begin{thebibliography}{99}

\bibitem{Planck2018}
Planck Collaboration (2020).
\textit{Planck 2018 results. VI. Cosmological parameters}.
A\&A 641, A6.
\url{https://arxiv.org/abs/1807.06209}

\bibitem{SH0ES}
Riess, A.G.\ et al.\ (2022).
\textit{A comprehensive measurement of the local value of the Hubble constant}.
ApJL 934, L7.
\url{https://arxiv.org/abs/2112.04510}

\bibitem{Freedman2019}
Freedman, W.L.\ et al.\ (2019).
\textit{Carnegie-Chicago Hubble Program VIII: $H_0$ from TRGB}.
ApJ 882, 34.
\url{https://arxiv.org/abs/1907.05922}

\bibitem{DESI2024}
DESI Collaboration (2024).
\textit{DESI 2024 VI: Cosmological constraints from BAO}.
\url{https://arxiv.org/abs/2404.03002}

\bibitem{KiDS1000}
Asgari, M.\ et al.\ (2021).
\textit{KiDS-1000 cosmology: cosmic shear constraints}.
A\&A 645, A104.
\url{https://arxiv.org/abs/2007.15633}

\bibitem{Bohme2025}
Böhme, M.\ et al.\ (2025).
\textit{Overdispersed radio source counts and excess radio dipole detection}.
\url{https://arxiv.org/abs/2509.16732}

\bibitem{Fontana2024}
Fontana, G.\ et al.\ (2024).
\textit{Pulsed electromagnetic field stimulation enhances neurite outgrowth
in neural cells and modulates inflammation in macrophages}.
ResearchGate: \url{https://doi.org/10.13140/RG.2.2.29234.17608}

\bibitem{EDGES2018}
Bowman, J.D.\ et al.\ (2018).
\textit{An absorption profile centred at 78 MHz in the sky-averaged spectrum}.
Nature 555, 67--70.
\url{https://doi.org/10.1038/nature25792}

\bibitem{HERAPhaseII}
HERA Collaboration (2025).
\textit{First results from HERA Phase II}.
\url{https://arxiv.org/abs/2511.21289}

\bibitem{EuclidOverview}
Euclid Collaboration (2024).
\textit{Euclid I: Overview of the Euclid mission}.
\url{https://arxiv.org/abs/2405.13491}

\bibitem{LIGODarkSiren}
LVK Collaboration (2024).
\textit{A dark standard siren measurement of the Hubble constant
following LIGO/Virgo/KAGRA O4a}.
MNRAS 535, 961.
\url{https://arxiv.org/abs/2404.16092}

\bibitem{XRISM2025}
XRISM Collaboration (2025).
\textit{Constraints on unidentified X-ray emission lines in stacked
spectrum of 10 galaxy clusters}.
\url{https://arxiv.org/abs/2510.24560}

\bibitem{Romer2026}
Römer, J.B.\ (2026).
\textit{The evolution of AFET: iterative analysis of UTAC dynamics,
biophysical resonances, and cosmological scaling laws} (v37).
Zenodo. \url{https://doi.org/10.5281/zenodo.18641306}

\end{thebibliography}

\end{document}
